\chapter{Literature}
In this chapter, we explore some classifications related to AI safety-critical systems. This is useful for providing a global overview of the context in which we will be moving.

%%%%%%%%%%%%%%%%%%%%%%%%%%%%%%%%%%%%%%%%%%%%%%%%
\section{Types of AI}
\label{sec:TAI}
In the literature, there are many Types of AI (TAI). In this section, we provide a brief recap of them, as they may be useful in the following sections.
\begin{itemize}
    \item \textbf{Connectionists} \\
    These are models where training is performed using an optimization technique, most commonly gradient descent. Examples include Neural Networks (e.g., CNNs, Transformers).
    \item \textbf{Bayesian} \\
        Results are obtained in a probabilistic manner.
    \item \textbf{Symbolists} \\
        These use logic to perform reasoning (e.g., decision trees and Prolog). Given that inference must validate all parts of the logic (i.e., each component of a boolean expression) learned during the training phase, these models are useful if we want an explainable model (e.g., in the \href{https://github.com/edoardosarri24/violation-prediction-in-realtime-system}{prediction of faults} in real-time scenarios).
    \item \textbf{Analogizers} \\
        Usually, these are classification models based on a similarity score (e.g., k-NN, SVM).
\end{itemize}

%%%%%%%%%%%%%%%%%%%%%%%%%%%%%%%%%%%%%%%%%%%%%%%%
\section{Automation Levels}
In the field of AI, we can define several levels of automation. A classification can be made, from the highest to the lowest, in the following way:
\begin{enumerate}
    \item \textbf{Autonomous systems} operate in an open environment. \\
        ISO 22989 defines this as the highest level of automation, in which the system can modify its behavior and goals without human control.
    \item \textbf{Heteronomous systems} operate in a (semi-)open environment. \\
        These include different levels of automation where requirements can be defined within a range but are not fully known, and human control can be variable.
    \item \textbf{Automatic systems} operate in a closed environment. \\
        In this case, the system executes in a predefined environment where requirements are known during the design phase.
    \item \textbf{Collaborative Robots} \\
        These are robots that work side by side with humans without having total independence.
\end{enumerate}

%%%%%%%%%%%%%%%%%%%%%%%%%%%%%%%%%%%%%%%%%%%%%%%%
\section{Safety Standards}
\label{sec:standard}
Moving from traditional methods of guaranteeing safety to using AI to achieve it requires defining a new class of standards. The most common problem is related to the probability of failure: the most critical Safety Integrity Levels (SIL, ASIL, or DAL) require a very low probability of failure (e.g., $10^{-9}$), whereas an excellent accuracy for an AI model is often 99\%. Thus, common AI models are not compliant with these requirements.
\begin{itemize}
    \item \textbf{Traditional safety standards} \\
        Their objective is to mitigate risk using programmable functions, known as Safety Instrumented Systems (SIS), which are embedded in the system. The quality of a SIS is defined by different measures depending on the domain: we have the Safety Integrity Level (SIL) for industry (IEC 61508), the Automotive Safety Integrity Level (ASIL) for the automotive sector (ISO 26262), and the Design Assurance Level (DAL) for aerospace (DO-178C).
    \item \textbf{AI Safety Standards} \\
        Having defined heteronomous systems (e.g., Advanced Driver Assistance Systems (ADAS)), we can see that failures can occur even in the absence of hardware or software faults. This means that these new algorithms must comply with the Safety of the Intended Functionality (SOTIF) standard (ISO 21448), bringing them into alignment with classical FuSa.
\end{itemize}

%%%%%%%%%%%%%%%%%%%%%%%%%%%%%%%%%%%%%%%%%%%%%%%%
\section{Points of Application}
\label{sec:application}
In this section, we recap the different types of AI applications. There are three categories that describe \textit{when} and \textit{on what} we can use an AI model.
\begin{itemize}
    \item \textbf{Product} \\
        A safety-critical system relies on one or more safety functions based on an AI model. This model is offline: training is performed before it is embedded in the system.
    \item \textbf{Runtime} \\
        This is similar to the previous category, but the model can update itself during execution.
    \item \textbf{Process} \\
        The system itself does not include an AI model, but it is designed and developed using AI models (e.g., generative AI to generate code or perform tests).
\end{itemize}

\myskip

Based on this classification, ISO 5469 defines three classes for when AI models can be included in a safety-critical system (i.e., product, runtime, and process) while remaining compliant with the standards defined in Section~\ref{sec:standard}.
\begin{itemize}
    \item \textbf{Class I} \\
        These systems are based on traditional AI models, such as those referred to as \textit{Symbolists} in Section~\ref{sec:TAI}. Since they are totally deterministic and formal verification can be performed, these models comply with classic FuSa.
    \item \textbf{Class II} \\
        The system contains an AI model that may not be compliant with FuSa on its own. Compliance is achieved through a run-time checker, often known as a \textit{safety bag}; if the model fails, the checker notices and makes a safe choice.
    \item \textbf{Class III} \\
        These are systems that cannot incorporate an AI model internally. In this case, the solution is usually to include the model, but engineers do not take responsibility for possible failures (e.g., ADAS in vehicles).
\end{itemize}